% ---------------------------------------------------------------------------
% Week 9 progress report for the supervisor. Compiles with pdflatex
% (Overleaf: New Project -> Upload Project, or paste into main.tex).
% Style follows 2026-10-01_progress_update.tex (cent sign via \cent).
% Sources: docs/planning/maker-simulator-preregistration.md (A12-A15),
% docs/planning/lumid-forward-test-preregistration.md (incl. L1),
% logs/v2_forecast/{smoke_test_2026-09,g2_W3_final_selection,tune_gbm_a15_summary}.log,
% runs/live/ (producer record).
% ---------------------------------------------------------------------------
\documentclass[11pt,a4paper]{article}

\usepackage[T1]{fontenc}
\usepackage[utf8]{inputenc}
\usepackage{textcomp}
\usepackage{amsmath}
\usepackage{booktabs}
\usepackage{tabularx}
\usepackage[margin=2.2cm]{geometry}
\usepackage{enumitem}
\usepackage{hyperref}

\newcommand{\cent}{\text{\textcent}}
\newcolumntype{L}{>{\raggedright\arraybackslash}X}
\setlist{itemsep=2pt,topsep=4pt}
\setlength{\emergencystretch}{3em}

\title{FYP Week 9 Progress Report:\\Learned Market-Making on Kalshi Event Contracts}
\date{Week 9 (12--16 October 2026)}
\author{}

\begin{document}
\maketitle
\vspace{-2em}

\section{Introduction}

This project asks whether a learned model can trade Kalshi event contracts profitably, and where it adds value
over simple rules. Since the last update (1 October), the learned market-making policy was carried from 2025 into
2026 conditions. Its 2025 version stopped working, a retrained version recovered most of the edge on validation,
and the strategy now publishes live signals to Lumid every hour. This report covers that work and asks for your
guidance on the questions in Section~\ref{sec:questions}, several of which decide how the forward test can be run.

\paragraph{Summary of progress.}
\begin{enumerate}
  \item \textbf{The 2025 policy (G) decayed in 2026.} On 7 September -- 4 October 2026 it returned $-1.41\%$, no
        better than quoting everything ($-1.45\%$). The weather markets it relied on changed character.
  \item \textbf{A retrained policy (G2) recovers the edge on validation.} Trained on Nov 2024 -- Aug 2026, it
        returned $+1.41\%$ on the same window, $+2.86$ points above quoting everything.
  \item \textbf{The edge holds month by month.} Walk-forward over five unseen months (Apr--Aug 2026), G2's model
        beat quoting everything in every month, by $+3.49$ points pooled (CI $+2.20$ to $+4.78$), with a pooled
        return of $+2.63\%$ (CI $+1.25$ to $+3.94\%$).
  \item \textbf{G2's results are not an artefact of its hyperparameters.} A pre-registered walk-forward sensitivity
        analysis (14 settings, 5 monthly folds) found every squared-error setting positive in every fold, while
        quoting everything was negative in every fold.
  \item \textbf{The live pipeline runs, but the Lumid forward test has not started.} The signal producer publishes
        hourly, but my deployed strategy has never run a decision cycle, and Lumid's tradable universe barely
        overlaps the markets the strategy trades.
\end{enumerate}

\section{Background}

\paragraph{Kalshi event contracts.} Each contract pays \$1 if an event occurs and \$0 otherwise, so its price is
a probability in cents. Many events are listed as \emph{mutually exclusive baskets}: several contracts (``legs''),
exactly one of which resolves YES, e.g.\ the outcomes of a football match, the bucket for a Fed decision, or the
range of tomorrow's high temperature. This project trades only such baskets.

\paragraph{Why mutually exclusive baskets.}
\begin{itemize}
  \item \textbf{They carry a hard constraint.} Because exactly one leg pays \$1, the fair prices of the legs must
        sum to \$1. Any departure is measurable at every moment, against a known target, without a model of the
        underlying event. This gave the project its original question: does a learned model add value over the
        constraint itself, or does the constraint suffice?
  \item \textbf{Each leg is priced in context.} The other legs of the same basket say something about each leg's
        fair price. A leg whose basket sums to more than \$1 is more likely overpriced. The policy uses this
        directly: four of its 14 features describe the basket (its size, how many of its legs are quoted, the sum
        of its mids, and the best other leg). The same structure defined the STGAT's graph, which linked each leg to its basket.
  \item \textbf{Outcomes are clean to score.} Each basket settles with exactly one YES, so every fill's P\&L at
        settlement is unambiguous. A basket that does not settle that way (e.g.\ a voided event) is easy to detect
        and exclude.
  \item \textbf{They cover many market types.} Sports results, economic releases, Fed decisions and weather all
        list as baskets. One method then applies across categories, and the results can be compared category by
        category.
\end{itemize}

\paragraph{Makers and takers.} A \emph{taker} crosses the spread and trades immediately; a \emph{maker} posts a
limit order and waits to be filled. Takers pay the spread and a fee on every trade. Makers earn the spread when
filled, but are exposed to \emph{adverse selection}: they tend to be filled precisely when an informed trader
knows the price is about to move against them.

\paragraph{Lumid.} The supervisor's platform (Lumid Studio) runs strategies written in a deterministic rule
language (LQT). A model cannot run inside a strategy, so model scores are computed outside, published to the
platform as named signals, and read by the strategy each cycle. The platform offers backtests on settled
contracts and a live paper-trading lane for forward tests.

\section{Motivation}

Earlier phases established three things that shape the current work:
\begin{enumerate}
  \item \textbf{Taker strategies on basket mispricing do not pay.} About 92.5\% of the apparent deviation of basket
        prices from \$1 is the bid--ask overround, not mispricing. At quoted prices, selling a basket at the bids
        nets $-0.07\cent$ before fees, against taker fees of $1.3$--$5.8\cent$ per basket.
  \item \textbf{The STGAT forecasts prices but does not improve trading.} It beats a shrinkage oracle on basket
        forecasting, but adding its forecasts to the trading policy changed test return by $-0.65$ points
        (CI $-1.40$ to $+0.09$).
  \item \textbf{A learned maker policy does pay, at least in 2025.} Policy G (a gradient-boosted model choosing
        which maker quotes to post) returned $+3.55\%$ net of fees on Aug--Nov 2025 (CI $+2.51$ to $+4.60\%$,
        34,086 fills), while posting every eligible quote lost $0.34\%$.
\end{enumerate}
The open questions this period were whether G's edge survives into 2026, how to evaluate it live on Lumid, and
whether its untuned hyperparameters matter.

\section{Method}

\subsection{Simulator}
A tape-driven replay of real Kalshi trades and hourly quotes. Each hour, candidate maker quotes are generated
for every eligible basket leg (join the touch, or improve by $1\cent$ when the spread is at least $3\cent$;
10 contracts; position limits per leg and per event). A resting quote is filled when later trades print through
it. Three queue-position models bound the outcome: \emph{optimistic}, \emph{pessimistic} (the pre-registered
verdict) and \emph{strict} (competing makers improve too). P\&L includes Kalshi's maker fee schedule and is
taken at settlement.

\subsection{Learned policy (G, G2)}
For each candidate quote, a gradient-boosted regression model (scikit-learn
\texttt{HistGradient\allowbreak{}Boosting\allowbreak{}Regressor}) predicts net P\&L per contract at settlement,
after fees, under pessimistic fills. The quote is posted only if the prediction exceeds $1\cent$. The model uses
14 features known at the quoting hour: side, price, price relative to mid, spread, mid, trades in the last 24h,
24h mid change, basket size, number of quoted legs, sum of basket mids, best other leg, hour, fee status and
category. Hyperparameters were fixed by judgement before any training (400 trees, learning rate 0.03, 31 leaves,
at least 300 samples per leaf, L2 penalty 1.0) and never searched on results.

\paragraph{G2.} The same design retrained on recent data. Three training windows were pre-registered and the
best one chosen on validation (Section~\ref{sec:setup}).

\subsection{Hyperparameter sensitivity (amendment A15)}
Fourteen settings of G2's own model, each changing one thing from G2 (leaf size, number of leaves, learning
rate, trees, L2 penalty, feature subsampling, time-ordered early stopping, recency weighting, absolute-error
loss, bagging), plus XGBoost with matched settings as a library benchmark. Each is trained walk-forward on
training months only and scored on the following month. The rule, fixed before scoring, selects the highest
\emph{worst-fold} return among settings keeping at least 200 fills in every fold. The aim is robustness, not peak
return.

\subsection{Live pipeline}
A producer on my machine runs hourly: it lists open Kalshi baskets, builds candidates and features with the
simulator's own code, scores them with each model, logs every score, and publishes the scores to Lumid as
signals. Replaying the validation window through this live path reproduced the offline scores exactly (0
mismatches in 2,142 scores). Six strategy arms are prepared for Lumid: G2 (primary), G2+ (an adverse-selection
variant), G (decay reference), two random-selection controls at the models' hourly share, and quote-everything.

\paragraph{Pre-registration.} Every model, window, threshold, selection rule and verdict was committed to the
repository before the corresponding result was computed, with dated amendments (A12--A15 for the policy, L1 for
the forward test).

\section{Experimental setup}\label{sec:setup}

\begin{center}
\small
\begin{tabularx}{\textwidth}{l L L}
\toprule
Study & Training & Evaluation \\
\midrule
G (2025) & Nov 2024 -- May 2025 & validation Jun--Jul 2025; \textbf{test Aug--Nov 2025} (scored once) \\
2026 check of G & (frozen G) & 7 Sep -- 4 Oct 2026 \\
G2 (2026 retrain) & W1: Jun--Aug 2026; W2: Mar--Aug 2026; W3: Nov 2024 -- Aug 2026 & \textbf{validation} 7 Sep -- 4 Oct 2026
  (window selection) \\
A15 sensitivity & every month before the fold's test month & 5 folds: Apr, May, Jun, Jul, Aug 2026 (training months only) \\
Forward test & (frozen G2-W3, G2+-W3, G) & live, from 6 Oct 2026; interim cut 28 Oct; end 11 Nov 2026 \\
\bottomrule
\end{tabularx}
\end{center}

\begin{itemize}
  \item \textbf{Metric.} Return = net P\&L after fees $/$ capital committed at fill prices, pessimistic fills.
        95\% confidence intervals by event-cluster bootstrap (fills in the same event are not independent).
  \item \textbf{Baselines.} Quoting every eligible candidate (no model), and random selection of the same share of
        candidates the model posts. Beating the random control shows that the model's \emph{choice} adds value,
        not just posting less.
  \item \textbf{Data.} Kalshi trades, hourly quotes and settlements: the historical archive to Nov 2025, and the
        Kalshi API for Dec 2025 onwards. The 23 basket series that qualified in 2025 are used for training.
\end{itemize}

\section{Results}

\subsection{The 2025 policy decayed in 2026}
Run unchanged on 7 September -- 4 October 2026, G returned $-1.41\%$ (CI $-4.10$ to $+1.36\%$; 10,599 fills,
329 events), the same as quoting everything ($-1.45\%$) and only slightly above random selection ($-1.95\%$).
The cause is a change in the weather markets, which made up 72\% of G's fills: 4.6 times more trade prints,
tighter spreads (the share of $2\cent$ spreads rose from 15\% to 45\% of candidates), and quote-everything
markouts turning negative ($+0.31$ to $-0.67\cent$ per contract). G's most confident decile, $+4.1\cent$ per
contract in 2025, was $-0.7\cent$ in 2026.

\subsection{Retrained G2 on validation}
\begin{center}
\small
\begin{tabularx}{\textwidth}{L r c r c r}
\toprule
Policy (validation, 7 Sep -- 4 Oct 2026) & Return & 95\% CI & \shortstack[r]{vs quote\\everything} & 95\% CI & Fills \\
\midrule
G (2025 model, unchanged) & $-1.41\%$ & $-4.10$ to $+1.36$ & $+0.04$ & $-2.35$ to $+2.63$ & 10,599 \\
G2-W1 (Jun--Aug 2026) & $+1.25\%$ & $-0.70$ to $+3.35$ & $+2.70$ & $+0.73$ to $+4.51$ & 9,293 \\
G2-W2 (Mar--Aug 2026) & $-0.36\%$ & $-2.76$ to $+2.26$ & $+1.09$ & $-1.40$ to $+3.38$ & 7,446 \\
\textbf{G2-W3 (Nov 2024 -- Aug 2026), selected} & $\mathbf{+1.41\%}$ & $-1.47$ to $+4.31$ & $\mathbf{+2.86}$ & $-0.23$ to $+5.69$ & 5,215 \\
Quote everything & $-1.45\%$ & $-2.40$ to $-0.53$ & -- & -- & 29,452 \\
\bottomrule
\end{tabularx}
\end{center}

\begin{itemize}
  \item \textbf{Retraining recovers the model's added value.} G2-W3 is $+2.86$ points above quoting everything and
        $+3.0$ points above random selection of the same share ($-1.60\%$). Its absolute return is positive but
        its CI includes 0: one month of validation is short.
  \item \textbf{In September, G2-W3's edge was outside weather.} Split by market group, G2-W3 returned $+6.73\%$ outside weather
        (CI $+2.4$ to $+11.4\%$; 1,725 fills) and $-1.60\%$ on weather (CI $-5.2$ to $+2.0\%$; 3,490 fills). On
        this, decided at group level and before any Lumid trading, weather series are excluded from the Lumid
        forward test (amendment L1). Section~\ref{sec:noise} qualifies this.
  \item \textbf{The three windows cannot be told apart.} Paired on the same events, W1 $-$ W2 is $+1.61$ points
        (CI $-1.06$ to $+4.20$), W3 $-$ W2 $+1.77$ (CI $-0.56$ to $+3.97$) and W3 $-$ W1 $+0.15$ (CI $-3.11$ to
        $+3.27$). W1 and W3 also earn in different places: W1, trained on recent months only, makes $+1.93\%$ on
        weather and loses outside it; W3, trained on two years, makes $+6.73\%$ outside weather and loses on it.
        A plausible reading is that recent data suits the changed weather market, while the smaller non-weather
        markets need a longer history; W2 sits between the two.
  \item The validation window was used to choose among three windows, so these figures are optimistic. The forward
        test is the first unseen data for G2.
\end{itemize}

\subsection{Hyperparameter sensitivity (A15)}
Offline return of the quotes each setting keeps (pessimistic fills), per walk-forward fold:

\begin{center}
\small
\setlength{\tabcolsep}{4pt}
\begin{tabular}{l r r r r r r r}
\toprule
Setting & Apr & May & Jun & Jul & Aug & \textbf{Worst} & Mean \\
\midrule
\textbf{G2 as is} & $+1.95$ & $+2.30$ & $+3.26$ & $+3.46$ & $+2.59$ & $\mathbf{+1.95}$ & $+2.71$ \\
\textbf{15 leaves (selected)} & $+2.61$ & $+3.08$ & $+2.71$ & $+3.60$ & $+2.56$ & $\mathbf{+2.56}$ & $+2.91$ \\
L2 penalty 10 & $+2.52$ & $+2.36$ & $+2.61$ & $+3.87$ & $+2.32$ & $+2.32$ & $+2.73$ \\
80\% of features per split & $+2.84$ & $+2.46$ & $+2.56$ & $+3.46$ & $+1.91$ & $+1.91$ & $+2.65$ \\
Other 9 squared-error settings & \multicolumn{5}{c}{all positive in every fold} & $+0.78$ to $+1.66$ & $+1.54$ to $+2.41$ \\
Absolute-error loss & $-0.51$ & $-0.07$ & $+0.48$ & $-0.47$ & $-0.29$ & $-0.51$ & $-0.17$ \\
XGBoost, matched (benchmark) & $+3.21$ & $+3.48$ & $+3.66$ & $+1.70$ & $+2.01$ & $+1.70$ & $+2.81$ \\
\midrule
Quote everything & $-0.73$ & $-0.79$ & $-0.26$ & $-1.07$ & $-1.67$ & & \\
\bottomrule
\end{tabular}
\end{center}

\begin{itemize}
  \item \textbf{The result does not depend on the settings.} Every squared-error setting is positive in every
        month; quoting everything is negative in every month.
  \item \textbf{The selected setting (15 leaves, simpler trees) is only modestly better:} $+0.61$ points on the
        worst fold, $+0.20$ on the mean, within the $\pm1$ point predicted before scoring. Choosing the best of 14
        flatters its margin, and it keeps fewer quotes (about 15\% of candidates vs 18--22\% for G2).
  \item \textbf{Changing library gains nothing:} XGBoost's mean is within $0.1$ points of G2's, with a lower worst
        fold.
  \item \textbf{Absolute-error loss fails}, as expected in hindsight: it predicts the median payoff, not the mean,
        so the $1\cent$ threshold keeps half of all candidates.
  \item These are offline scores without inventory interaction. The pre-registered next step re-runs the selected
        setting and G2 in the full simulator on the same months.
\end{itemize}

\subsection{Is G2's edge real or noise?}\label{sec:noise}
To judge this beyond one validation month, G2's exact model (A12 settings) was scored on each of the five A15
walk-forward months, with event-clustered 95\% CIs added after the pre-registered analysis:

\begin{center}
\small
\begin{tabular}{l r c r c r}
\toprule
Test month & \shortstack[r]{Return of\\quotes kept} & 95\% CI & \shortstack[r]{Minus quote\\everything} & 95\% CI & Events \\
\midrule
Apr 2026 & $+1.95\%$ & $-0.36$ to $+4.23$ & $+2.68$ & $+0.40$ to $+4.97$ & 496 \\
May 2026 & $+2.30\%$ & $-0.43$ to $+5.02$ & $+3.09$ & $+0.48$ to $+5.68$ & 496 \\
Jun 2026 & $+3.26\%$ & $-0.38$ to $+6.84$ & $+3.52$ & $-0.18$ to $+7.12$ & 303 \\
Jul 2026 & $+3.46\%$ & $+0.09$ to $+6.90$ & $+4.53$ & $+1.42$ to $+7.76$ & 344 \\
Aug 2026 & $+2.59\%$ & $-0.90$ to $+5.89$ & $+4.25$ & $+0.90$ to $+7.62$ & 418 \\
\midrule
\textbf{Pooled} & $\mathbf{+2.63\%}$ & $+1.25$ to $+3.94$ & $\mathbf{+3.49}$ & $+2.20$ to $+4.78$ & 2,012 \\
\bottomrule
\end{tabular}
\end{center}

\begin{itemize}
  \item \textbf{Real: the model's choice of quotes adds value.} It beats quoting everything in all five months
        (four of five CIs exclude 0), as it did on the 2025 test and the September validation.
  \item \textbf{Real but modest: absolute profit,} about $+2.6\%$ pooled. In a single month the CI usually
        includes 0 even though the pooled CI does not.
  \item \textbf{Noise: the ranking of training windows and single-month group figures.} Over these five months,
        weather quotes returned $+2.34\%$ (CI $+0.61$ to $+4.05$) and non-weather $+3.12\%$ (CI $+1.22$ to
        $+5.01$). September's weather loss for W3, on which the weather exclusion (L1) was decided, may therefore
        be one poor month rather than a lasting change. L1 stands as pre-registered, but this is reported.
  \item \textbf{Implication for the forward test.} A few weeks, at the current coverage, will most likely give a CI
        that includes 0 even if the edge is real. The paired comparison against quoting everything and random
        selection is the claim the forward test can best support.
  \item These are offline scores (no inventory limits, pessimistic fills). Under the strict fill model, where
        competing makers also improve, most of the edge disappeared in earlier tests.
\end{itemize}

\subsection{Live pipeline and Lumid status}
\begin{itemize}
  \item \textbf{Producer.} Running since 6 October. A full cycle covers about 1,900 open baskets and 12,800 legs
        and takes about 32 minutes. Signal publishing to Lumid succeeds (every batch accepted in the latest hours),
        after a Lumid ingest outage from 9 October 11:20 UTC to 10 October 07:00 UTC.
  \item \textbf{Coverage is low.} Only 25 of the roughly 110 hours between 6 October 16:00 UTC and 11 October 05:00
        UTC were logged, mainly because my machine and home connection are off overnight. Missed hours cannot be
        scored afterwards, and the gaps fall at the same times each day.
  \item \textbf{The strategy has never run on Lumid.} \texttt{g2\_maker\_v1} (registered 7 October, region
        \texttt{eu-west-1}) shows \texttt{live\_enabled = false}, and its decision-cycle feed is empty.
  \item \textbf{The universe barely overlaps.} On 10 October, Lumid's Kalshi tradable list held 84 instruments,
        mostly crypto price-threshold contracts. Of G2's 13,358 candidate quotes that hour, only 176 were in series
        on the list; at instrument level the overlap was a few Fed-decision legs. Some listed sports markets were
        stale (finished games still marked active, current fixtures missing).
  \item Your reply that paper fills are instant at the limit price, with no P\&L computed, means the paper lane
        cannot measure maker fills as specified in amendment L1 (Section~\ref{sec:questions}, Q1--Q2).
\end{itemize}

\section{Limitations}
\begin{itemize}
  \item One month of 2026 validation, already used for model selection; wide CIs.
  \item The simulator's fill rule is an assumption about queue position; the strict model shows the edge
        largely disappears if competing makers also improve their quotes.
  \item The A15 scores ignore inventory limits; the full-simulator re-check is pending.
  \item Forward-test coverage is currently limited by my machine's uptime, and a few weeks of data have little
        power: a real edge of the size seen above would often not reach significance.
  \item The weather exclusion rests on a single month, which the five-month evidence does not support as strongly.
\end{itemize}

\section{Next steps}
\begin{enumerate}
  \item \textbf{Finalise the experimental setup,} based on your answers to Q1--Q6, and record it as a dated
        amendment before any live result is seen:
        \begin{itemize}
          \item \emph{Backtesting:} which historical windows are reported (2025 held-out test, 2026 walk-forward,
                or both), how P\&L is computed (my simulator on recorded tape, settlement and fees), and the
                standard a result must meet (test length, significance).
          \item \emph{Forward testing:} the forward-test period and its interim cut, how its P\&L is measured
                (forward-period replay of logged scores, the Lumid paper lane, or both), and the start date if
                Lumid's universe is widened.
        \end{itemize}
  \item \textbf{Fix the required baselines:} the strategy baselines (quoting everything, random selection of the
        same share, the rule-based maker, the 2025 model), the model baselines (XGBoost, STGAT features, and a
        linear model if wanted), and the financial benchmarks (S\&P 500 and T-bills: Sharpe ratio, maximum
        drawdown, correlation), and which periods each covers.
  \item \textbf{Decide which model executes the strategy:} keep G2 (A12 settings) as the primary model, or add the
        tuned setting (15 leaves) as an extra arm after a full-simulator re-check against G2 on the same months.
  \item Improve producer uptime (a long-running host instead of my laptop).
  \item Interim forward-test cut on 28 October; final report due 3 November.
\end{enumerate}

\section{Questions for you}\label{sec:questions}

\subsection*{Evaluation}
Your reply on fills explained that the paper lane fills a resting limit order immediately at its limit price, with
no queue and no fees, and computes no P\&L for user strategies, and that the more reliable route to an edge
estimate is a backtest on recorded Kalshi tape with resolved settlement and Kalshi's fee schedule. My simulator does
exactly this, so I would like to confirm how that applies to the FYP evaluation.
\begin{enumerate}
  \item \textbf{Is P\&L evaluated by backtest?} Should the P\&L I report come from backtests on recorded Kalshi tape
        using my own simulator (queue-aware fills, settlement and fees), rather than from Lumid's paper lane? This
        would apply to two kinds of backtest:
        \begin{itemize}
          \item \emph{historical:} held-out months the model never trained on (e.g.\ the 2025 test);
          \item \emph{forward-period:} the live weeks, replayed through the same simulator using the scores my
                producer logged each hour before outcomes were known.
        \end{itemize}
        If a queue-aware paper fill model becomes available on Lumid before the report, should its P\&L replace
        the backtest or be reported beside it?
  \item \textbf{Purpose of the Lumid forward test.} Given that the paper lane records decisions only, is the forward
        test on Lumid mainly to show that the strategy runs live on the platform, with its P\&L measured by the
        forward-period backtest? Or is performance on the paper lane itself part of the evaluation?
  \item \textbf{Is a historical backtest required, and on which windows?} I propose three layers, each strictly
        forward in time:
        \begin{itemize}
          \item \emph{2025 held-out test:} train Nov 2024 -- May 2025, validation Jun--Jul 2025, test Aug--Nov 2025
                (the 7/2/4 split you approved on 28 September; scored once);
          \item \emph{2026 walk-forward:} five monthly folds, Apr--Aug 2026, each trained only on earlier months.
                This tests G2's training recipe rather than the exact deployed model, whose training data include
                these months;
          \item \emph{forward period:} the deployed G2 (trained to Aug 2026, validated 7 Sep -- 4 Oct 2026) on the
                live weeks to 11 November (interim cut 28 October). This is the only data G2 itself has not seen.
        \end{itemize}
        Is this appropriate, or would you prefer one long backtest on a single split? Is a historical backtest
        needed at all if the forward-period backtest is the main result? Changing the windows is cheapest now; some
        test months have already been scored once, and reusing them in a new split would have to be disclosed.
  \item \textbf{Standard for the backtest.} Is there a minimum test length or number of events? Must the profit be
        statistically significant (95\% CI excluding 0), or is a positive return reported with its CI acceptable?
        A few weeks of forward data will most likely give a CI that includes 0 even if the edge is real.
  \item \textbf{Which baselines?} I currently compare against:
        \begin{itemize}
          \item \emph{strategy baselines} on the same markets and hours: quoting every eligible order (no model),
                random selection of the same share of quotes, a hand-set rule-based maker, and the 2025 model;
          \item \emph{model baselines}: XGBoost, and the policy with STGAT features added.
        \end{itemize}
        For the financial benchmarks you mentioned, I plan to compare risk-adjusted returns (Sharpe ratio, maximum
        drawdown, correlation) with the S\&P 500 and 3-month T-bills over the same dates. Should this comparison
        cover the historical backtests as well as the forward period? Are there other baselines you expect, for
        example a simple linear model, or a buy-and-hold position in the same contracts?
  \item \textbf{Tuned or original model?} The sensitivity analysis selected a simpler setting (15 leaves) that is
        $+0.61$ points better on the worst fold. Is it worthwhile to add it as an extra forward-test arm, or should
        I keep the original G2 as the only model and report the analysis as a robustness check? Adding it would
        give it a shorter live window than the other arms.
\end{enumerate}

\subsection*{The Lumid universe and platform}
\begin{enumerate}[resume]
  \item \textbf{The live universe.} The Kalshi universe on Lumid is too incomplete for my strategy to trade.
        It trades mutually exclusive baskets (e.g.\ football match outcomes, Fed-decision buckets), but the tradable
        list on my account held 84 instruments on 10 October, mostly crypto price-threshold contracts. Only 176 of
        13,358 candidate quotes in one hour were in series on that list. Some listed markets are also stale: the
        only MLS games were finished matches from 1 and 6 October still marked active, while current fixtures were
        missing. Could the series my strategy trades be included in the universe refresh for my account? If that
        is not possible, how would you like me to run the forward test?
  \item \textbf{My strategy has never run on the paper lane.} \texttt{g2\_maker\_v1} (strategy id
        \texttt{af1a4954-e659-4c18-b9a9-93c999719dd8}, registered 7 October 21:01 UTC, region
        \texttt{eu-west-1}) shows \texttt{live\_enabled = false} with no \texttt{live\_enabled\_at}, and its
        decision-cycle feed is empty. The documentation says a strategy forward-tests from deployment, with no
        start button. Could the operator check whether it needs enabling, or whether I have missed a step?
  \item \textbf{Forward-test attribution.} The platform's evaluator reports that scorecards carry no
        \texttt{strategy\_id} because the build predates attribution, and that the fleet needs redeploying. Could
        the operator redeploy it, so that decisions can be attributed to my strategies?
  \item \textbf{Confirming that signals arrive.} My producer's batches are accepted (HTTP 200), but I cannot query
        the signal history to confirm that \texttt{g2\_bid} / \texttt{g2\_ask} rows are stored. Could the operator
        confirm that they are?
  \item \textbf{A long-running host for the producer.} Could I have a sandbox without the 24-hour expiry to host my
        signal producer? My home connection is cut overnight, so only about a quarter of hours have been logged
        so far, and the gaps fall at the same times each day.
\end{enumerate}

\end{document}
